% SPDX-License-Identifier: GPL-3.0-or-later
% Copyright (C) 2026 Alessandro Greco
%% Русское альтернативное описание EPUB.
\altfig{Схема обмена на TLS-канале клиент — сторожевой узел. После TLS обе стороны отправляют VERSIONS с CircID 0; одна стрелка сторожевого узла объединяет четыре отдельные ячейки VERSIONS, CERTS, AUTH\_CHALLENGE и NETINFO. Клиент отправляет NETINFO, создаёт первый переход с CREATE2/CREATED2 и расширяет цепь до промежуточного и выходного узлов двумя RELAY\_EARLY [EXTEND2] и ответами RELAY [EXTENDED2]. BEGIN показан в RELAY\_EARLY лишь как допустимый пример, не обязательное правило; затем следуют CONNECTED, запрос и ответ DATA, при необходимости SENDME и завершение DESTROY. Сторожевой узел наблюдает ячейки своего канала, но не читает сообщения для следующих узлов. Названия в скобках поясняют логические команды.}
% Последовательность ячеек на канале клиент — сторожевой узел.
\begin{tikzpicture}[font=\sffamily\small,>=Stealth]
  \node[client] (c) at (0,0) {Клиент};
  \node[relay] (g) at (7,0) {Сторожевой};
  \draw[torgrigio] (c.south) -- ++(0,-12.2);
  \draw[torgrigio] (g.south) -- ++(0,-12.2);
  \newcommand{\av}[3]{\draw[cellaf] (0,#1) -- node[above,font=\ttfamily\footnotesize]{#2} (7,#1);
    \node[font=\scriptsize,text=torgrigio,anchor=west] at (7.4,#1) {#3};}
  \newcommand{\ind}[3]{\draw[risposta] (7,#1) -- node[above,font=\ttfamily\footnotesize]{#2} (0,#1);
    \node[font=\scriptsize,text=torgrigio,anchor=west] at (7.4,#1) {#3};}
  \av{-1.1}{VERSIONS}{CircID 0}
  \ind{-1.8}{VERSIONS, CERTS, AUTH\_CHALLENGE, NETINFO}{CircID 0}
  \av{-2.5}{NETINFO}{канал готов}
  \av{-3.5}{CREATE2}{CircID 0x8\ldots, 1-й переход}
  \ind{-4.2}{CREATED2}{}
  \av{-5.2}{RELAY\_EARLY [EXTEND2]}{к промежуточному}
  \ind{-5.9}{RELAY [EXTENDED2]}{}
  \av{-6.9}{RELAY\_EARLY [EXTEND2]}{к выходному}
  \ind{-7.6}{RELAY [EXTENDED2]}{}
  \av{-8.6}{RELAY\_EARLY [BEGIN]}{поток к сайту}
  \ind{-9.3}{RELAY [CONNECTED]}{}
  \av{-10.1}{RELAY [DATA]}{HTTP-запрос}
  \ind{-10.8}{RELAY [DATA] \ldots}{ответ, SENDME}
  \av{-11.7}{DESTROY}{цепь завершается}
\end{tikzpicture}
